\begin{abstract}
This thesis involves modelling, parameter estimation and control strategies for a tendon-driven hand exoskeleton. The system is intended to assist hand function in spinal-cord-injured patients with the goal of controlling a \mbox{6-DOF} exoskeleton. The exoskeleton is actuated by brushed DC motors, pulling on a cable connected to a glove mounted on the hand of the user. The cable runs through a bowden tube, causing friction which the system has to overcome. In addition, varying stiffness of the user's hand and limited sensing at the hand introduce uncertain dynamics, and unobservable states, which the work in this thesis attempts to address. The thesis focuses on a single finger as a representative subsystem of the full 6-DOF exoskeleton.

A planar three-link finger model was derived using the Euler-Lagrange equations, including joint stiffness and damping. A brushed DC motor model was also presented, and relevant motor parameters were estimated from datasheet values for the relevant motor. For the parameters not available from the datasheet, a gradient-based parameter identification method was developed, and evaluated in simulation. For connecting the motor model to the finger model, a cable constraint was introduced, relating the angle of the motor to the angles of the finger joints, and the length of the cable.

A discrete-time current controller was designed and stability was guaranteed by analyzing the discrete-time closed-loop system. Simulations indicated that the controller could track the armature current rapidly compared to the mechanical dynamics, supporting the use of armature current as a control input for torque and position control. For force control, integrated fingertip force sensors were considered, showing promising results in reaching the target force, with friction from the bowden tubes present. With only motor-side sensing, the applied force would get a steady-state error. For position control, a model reference adaptive control scheme was designed using the motor angle as a proxy for fingertip position, given the unobservable states of the finger. The adaptive controller improved tracking over time, but the adaptive gains were observed to continue growing, likely due to nonlinearities, saturation, and deviations in the model compared to the model assumptions. The results therefore suggested that the adaptive scheme is most suitable as an automatic tuning method for finding fixed controller gains, simplifying the tuning of the controller for a new user.

Overall, the proposed controllers provide a simulation-based proof of concept for future implementation on the physical exoskeleton. Experimental validation is required to assess the performance of the controllers in practice, and to further refine the models and control strategies based on real-world data. 
\end{abstract} 